\abstract{
In this thesis, we will address the problem of capturing a tumbling, i.e., noncooperative, object in orbit. A tumbling object, also known as the target, can be a malfunctioning satellite or space debris that must be stabilized or removed to prevent collisions with other surrounding assets. Typical approaches use agents, called chasers, that attach to the tumbling object to stabilize or change the target’s orbit. While rendezvous and docking with cooperative objects has been possible since Project Gemini in the 1960s, effective strategies for tumbling objects that account for uncertainty in the target’s orbit have yet to be developed. We propose a novel adaptive tube model predictive control (MPC) formulation that is composed of a rendezvous (initial approach) and docking phase with the target. Importantly, the formulation includes uncertainty in the eccentricity and the drag parameters of the target's orbit. The approach is able to ensure all safety constraints are satisfied while also being able to incorporate online data to improve the prediction model, resulting in less conservative behavior. Simulation results of a single chaser-target scenario illustrate the approach. Moreover, an optimization-based method is proposed for computing the optimal amount of thrust to stop the tumble of the target once the chaser has docked with the target. A multi-chaser strategy is also explored.
}